\documentclass[11pt,a4paper]{article}
\usepackage[margin=24mm]{geometry}
\usepackage{amsmath,amssymb,booktabs,microtype}
\usepackage[hidelinks]{hyperref}
\newcommand{\dd}{\mathrm d}
\newcommand{\cN}{\mathcal N}
\title{Route G2-H: A Physical Source Budget and a Decisive Yield Bound}
\author{On-shell LHC Higgs production, not an invented dark beam}
\date{23 September 2026}
\begin{document}
\maketitle
\begin{abstract}
We select an existing accelerator process, Higgs production at the LHC,
and the same conditional $X$-only portal as G2-X. Production and detection
are tied to one coupling. A boost-independent, single-pass bound shows
that even perfect interception of the Run-2 dark-particle budget cannot
give a useful elastic xenon recoil sample in the retained energy window.
This closes one physical realization without a detailed beam simulation.
The direct Higgs source also erases the old sign preparation: opposite
compact labels have identical timing, direction and momentum likelihoods.
Neither result rules out the entire theory or constructs the original
driven mode-conversion source.
\end{abstract}

\section{Actual inputs and a different source preparation}
Keep G2-X's tree EFT and mass matching:
\begin{equation}
 \Delta\mathcal L=-\lambda_X H^\dagger H\sum_j|x_j|^2,\qquad
 \mu_j=\Lambda\sqrt{25+j^2},\qquad
 0\leq\lambda_X\leq\min(1,50\Lambda^2/v^2).
\end{equation}
The scale domain remains $1\leq\Lambda\leq30$ GeV. Use $m_h=125.08$ GeV,
$v=246$ GeV, and $\Gamma_{\rm SM}=.00410$ GeV as in G2-X. This branch
uses the \emph{drive-off} vacuum and ordinary Higgs decay, not a hidden
implementation of $g(t)$. It neither supplies a prepared $\Phi$ beam
nor performs the original $\Phi X$ conversion. The driven G2 population
and its priors remain a separate, unrealized source hypothesis.

Choose one LHC interaction point's published 13 TeV Run-2 luminosity
$\mathcal L_{\rm int}=139$ fb$^{-1}$. The ATLAS production paper quotes
the SM expectation $\sigma_h=55.6\pm2.5$ pb and an agreeing measurement
\cite{production}. We use the \emph{SM expectation}, consistent with the
SM-production premise of the invisible-Higgs bound; the visible-channel
measurement itself assumes SM branching fractions and is not an
independent production input after changing those fractions.
The small Higgs-mass convention difference in that production calculation
is not resolved beyond its quoted precision.
\begin{equation}\label{eq:budget}
 N_h=\mathcal L_{\rm int}\sigma_h=7.7284\,10^6,\qquad
 N_{X+\bar X}\leq2N_h B_{\rm inv}^{\max}=1.6538776\,10^6,
 \quad B_{\rm inv}^{\max}=.107 .
\end{equation}
The latter is the conditional observed ATLAS allowance at 95\% CL
\cite{invisible}, not a detected dark-particle sample. It assumes escaping
invisible modes and SM production/visible widths. No combined confidence
level is assigned to our subsequent theory/geometry bound.

The target column is based on XENON1T's approximately 2 tonne active
cylinder, radius 47.9 cm, height 97 cm \cite{xenon}. The detector is at
Gran Sasso, not at the LHC; no installation, transport link or concurrent
data collection is asserted. Perfect interception below is a deliberately
generous mathematical envelope. Rounded geometry is sufficient to test
an orders-of-magnitude shortfall.

\section{The source spectrum and flux functional}
For the diagonal vertex $-i\lambda_Xv$, two-body phase space gives
\begin{align}
 \beta_j&=\sqrt{[1-4\mu_j^2/m_h^2]_+},\qquad
 \Gamma_j=\frac{\lambda_X^2v^2}{16\pi m_h}\beta_j,\\
 B_j&=\frac{\Gamma_j}{\Gamma_{\rm SM}+\sum_k\Gamma_k},\qquad
 p_j^*=\frac{m_h\beta_j}{2},\quad E_j^*=\frac{m_h}{2}.
\end{align}
Kinematic statements below concern open modes only.
Each signed label is a distinct complex field and its decay creates
one particle and one antiparticle. Their compact momenta are $j\Lambda$
and $-j\Lambda$; a uniform Higgs conserves their sum. Do not put another
factor of two in $\Gamma_j$ or reuse the old selected-collision weights.
The twofold particle yield is already explicit in \eqref{eq:budget}.

Let $\dd F_h(P)$ be a normalized physical Higgs four-momentum distribution,
and let $L_P$ boost its rest frame to the laboratory. The single-particle
production measure is
\begin{equation}\label{eq:source}
 \frac{\dd N_{X+\bar X}}{\dd^3p}=
 N_h\sum_j B_j\int\dd F_h(P)\int\frac{\dd\Omega_*}{4\pi}
 \sum_{\epsilon=\pm1}
 \delta^3\!\left(\mathbf p-
    [L_P(E_j^*,\epsilon p_j^*\hat{\mathbf n}_*)]_{\rm spatial}\right).
\end{equation}
It integrates to $2N_h\sum_jB_j$. A real LHC source is boosted;
setting every Higgs at rest is not a physical flux prediction. Published
marginals alone are not silently promoted to a full correlated
source, transport and timing sample.

For a point source at distance $d$, a surface element subtending
$\dd\Omega=(\cos\theta/d^2)\dd A$ has fluence
\begin{equation}
 \frac{\dd\mathcal F_j}{\dd E\,\dd A}
 =\frac{\cos\theta}{d^2}\,
   \frac{\dd N_j}{\dd E\,\dd\Omega}\,T_j(E,\Omega),
 \qquad 0\leq T_j\leq1,
\end{equation}
under straight, nonmultiplying transport. Fluence is cm$^{-2}$,
not cm$^{-2}$s$^{-1}$. A flux additionally needs the actual
luminosity time profile, travel time and detector live time.
There is no reason to divide the Run-2 total by XENON's 180.7-day
search interval or to multiply by its kg-day exposure again.

More generally an incident particle with momentum $p$ on path $\gamma$
contributes, in the dilute elastic approximation,
\begin{equation}\label{eq:path}
 P_{\rm acc}=\int_\gamma \dd\ell\, n_{\rm Xe}
 \sum_A\xi_A\int_{E_0}^{E_1}\dd E_R\,
 \frac{\dd\sigma_A}{\dd E_R}
 \mathcal R_{\rm ROI}(E_R,\mathbf x,t),
 \quad (E_0,E_1)=(.7,50)\ {\rm keV}.
\end{equation}
The released S2 matrix supplies a \emph{volume-averaged} response, not
an arbitrary narrow beam's spatial response. Our decisive bound uses
$\mathcal R_{\rm ROI}\leq1$ everywhere, so it does not rely on applying
that average where invalid. The code separately checks the actual
response fold, including moving recoil endpoints, as a unit-column
kernel validation. No Gaussian replacement is introduced.

\section{A bound independent of boosts, geometry acceptance and fine tuning}
The retained G2-X scalar-current kernel is
\begin{align}
 \frac{\dd\sigma_A}{\dd E_R}
 &=\frac{\lambda_X^2 f_N^2m_N^2 A^2m_A}
 {8\pi p^2(m_h^2+2m_AE_R)^2}\,
 F_A^2(\sqrt{2m_AE_R})\,
 \mathbf1_{0<E_R<E_{\max}},\\
 E_{\max}&=\frac{2m_Ap^2}{s},\qquad
 s=\mu_j^2+m_A^2+2m_A\sqrt{\mu_j^2+p^2}.
\end{align}
Here $f_N=.308$, $m_N=.939$ GeV; natural isotope fractions and the
Helm conventions are unchanged from G2-X. Since $F_A^2\leq1$,
$\mathcal R\leq1$, the propagator denominator is at least $m_h^4$,
and $s\geq(\mu_j+m_A)^2$, integration over the \emph{admitted window}
gives
\begin{align}
 \sigma_{A,\rm acc}
 &\leq \frac{\lambda_X^2 f_N^2m_N^2 A^2m_A}
 {8\pi p^2m_h^4}\,E_{\max}\\
 &\leq\frac{\lambda_X^2 f_N^2m_N^2 A^2m_A^2}
 {4\pi m_h^4(\mu_j+m_A)^2}
 \equiv U_A(\mu_j,\lambda_X). \label{eq:majorant}
\end{align}
Using $E_{\max}$ in this inequality extends a mathematical \emph{constant
majorant}, not a nuclear model or calibration outside $[E_0,E_1]$.
In particular it does not claim a total coherent cross section at
MeV/GeV recoils. The $p=0$ accepted count is zero because $E_0>0$.

With uniform xenon density inferred from the rounded active volume,
\begin{equation}
 \ell_{\max}=\sqrt{(2r)^2+h^2}=136.333\ {\rm cm},\qquad
 \mathcal C_{\max}=\frac{M}{\pi r^2h}\ell_{\max}
                  \frac{N_A}{\bar M}
 =1.78874\,10^{24}\ {\rm nuclei\,cm^{-2}}.
\end{equation}
Any single straight traversal has at most this column. Thus even if
every produced particle is assigned the maximum chord,
\begin{equation}\label{eq:bound}
 N_{\rm acc}\leq
 2N_h\mathcal C_{\max}\sum_jB_j\sum_A\xi_A U_A(\mu_j,\lambda_X).
\end{equation}
This is an upper bound for the retained elastic recoil component,
not a simulated yield or a background likelihood. The inferred opacity
is tiny, consistent with the single-scattering approximation.

For a uniform bound across \emph{all} scales, any on-shell source has
$\Lambda<m_h/10$ and $\mu_j\geq5\Lambda$. The retained vacuum rule gives
$\lambda_X^2\leq2500\Lambda^4/v^4$. Every term is bounded using
\begin{equation}
 \frac{\dd}{\dd\Lambda}
 \log\frac{\Lambda^4}{(5\Lambda+m_A)^2}
 =\frac4\Lambda-\frac{10}{5\Lambda+m_A}>0.
\end{equation}
Taking the upper endpoint yields
\begin{align}
 \lambda_{\rm sup}&=\frac{m_h^2}{2v^2}=.1292634,\\
 U_{\rm sup}&=\sum_A\xi_A U_A(m_h/2,\lambda_{\rm sup})
             =1.33848\,10^{-36}\ {\rm cm}^2 .
\end{align}
\begin{equation}
 \boxed{N_{\rm acc}\leq2N_h B_{\rm inv}^{\max}
       \mathcal C_{\max}U_{\rm sup}=3.95969\,10^{-6}.}
 \label{eq:global}
\end{equation}
At $\Lambda\geq m_h/10$ the on-shell source is closed. Hence the bound
covers the full 1--30 GeV domain. This envelope does not require its
separate factors to be simultaneously attainable: exactly at the
threshold the actual production is zero. No search for a specially
tuned threshold point is needed.

An explicitly generous stress card, $\sigma_h=60$ pb, $f_N=.326$ and
$\mathcal C_{\max}=2\,10^{24}$ cm$^{-2}$, gives $5.35247\,10^{-6}$.
Even replacing $B_{\rm inv}^{\max}$ by one gives only
$3.70064\,10^{-5}$ under the other nominal premises. These are stress
envelopes, not confidence intervals. At least $7.58\,10^5$ times the
nominal source exposure would be needed for three events even under
\eqref{eq:global}'s impossible perfect interception; no extrapolated
accelerator performance is claimed.

At fixed masses, $\Gamma_j=\lambda_X^2u_j$ and
$\sigma_{\rm acc}=\lambda_X^2s_j$, so with $x=\lambda_X^2$ a
source-linked yield obeys
\begin{equation}
 N_{\rm acc}=\frac{A\lambda_X^4}{\Gamma_{\rm SM}+U\lambda_X^2},
 \quad A,U\geq0,\qquad
 \frac{\dd}{\dd x}\frac{x^2}{\Gamma_{\rm SM}+Ux}
 =\frac{x(2\Gamma_{\rm SM}+Ux)}{(\Gamma_{\rm SM}+Ux)^2}\geq0 .
\end{equation}
Thus choosing the allowed coupling ceiling is already optimistic;
the old independent-flux feasibility map must not be optimized by
choosing an unrelated large source luminosity.

\section{Mode information changes with the source}
For direct tree-level Higgs production, $\mu_j=\mu_{-j}$,
$\Gamma_j=\Gamma_{-j}$, and both scalar decays are isotropic. For any
common Higgs momentum distribution, the conditional laboratory
four-momentum measures in \eqref{eq:source} are identical for $j$ and
$-j$. The uniform, diagonal probe has identical matrix elements too.
For any measurement $z$ constructed from four-dimensional momentum,
direction, flight time and the same scattering response,
\begin{equation}
 P(z\mid j=+1)=P(z\mid j=-1),\qquad
 P(j=\pm1)=\tfrac12,\qquad I(j;z)=0,\qquad e_{\rm Bayes}=\tfrac12 .
\end{equation}
For example $t=d/(c\beta)$, $\beta=p/\sqrt{p^2+\mu_j^2}$, has the
same distribution for both signs. An independent clock or better
momentum instrument cannot distinguish equal likelihoods. This is
a statement about this production/probe channel, not an assertion of
an exact symmetry of the full $\Phi$-holonomy theory at all loops.
Different $|j|$ can have different speeds and thresholds; they are
not the same problem as resolving the sign of a degenerate pair.
The old asymmetric collision sample's 14.057\% S2 error is therefore
not transferred to this new source.

The selected $|j|=1$ source closes already at
$\Lambda=m_h/(2\sqrt{26})=12.2651$ GeV, before the last $j=0$ threshold.
This exposes a misleading intuition from a coupling-only map:
the high-scale region with less restrictive on-shell Higgs width
also lacks this on-shell source.

\section{Decision, next choices, and boundaries}
This stage rejects the Run-2 on-shell-Higgs plus single-pass xenon
\emph{retained-window elastic} chain as a useful event-count realization.
The verdict does not depend on a fitted Gaussian, a chosen boost
distribution or an achievable dark-particle focusing system.
The verifier checks source normalization, every open signed mode,
recoil-bin clipping using the unchanged public response, the continuous
bound, and Lorentz/sign identities. Its checks validate implementation,
not empirical existence of $X$.

Do not now build a detailed transport simulation for this already
insufficient chain. Two distinct next questions are more productive:
\begin{enumerate}
 \item Use a collider production/missing-momentum channel to avoid
 a second tiny scattering probability. On-shell Higgs invisibility
 constrains an inclusive width, not a resolved tower or its signs.
 Off-shell thresholds or a tagged missing invariant mass could test
 masses, but need a separately specified source and analysis.
 \item If signed-mode conversion is the indispensable claim, supply
 a genuinely mode-sensitive preparation or coupling. The original
 $\Phi$ collision correlations could provide it only after a physical
 $\Phi$ source is built; changing the detector resolution does not.
\end{enumerate}
Off-shell Higgs production, astrophysical populations, other targets,
electronic/inelastic channels, energies outside the admitted response,
trapping/recirculation and additional interactions are not ruled out
by the on-shell budget. None is silently added. The GeV-scale
classical pump remains unimplemented, and Route F remains independent.

\begin{thebibliography}{9}
\bibitem{production} ATLAS Collaboration, JHEP 05 (2023) 028,
\href{https://arxiv.org/html/2207.08615v2}{arXiv:2207.08615v2},
Secs.~1--2 and 5.
\bibitem{invisible} ATLAS Collaboration, Phys. Lett. B 842 (2023) 137963,
\href{https://arxiv.org/abs/2301.10731v2}{arXiv:2301.10731v2}.
\bibitem{xenon} XENON Collaboration, Phys. Rev. Lett. 123 (2019) 251801,
\href{https://arxiv.org/html/1907.11485v2}{arXiv:1907.11485v2};
\href{https://github.com/XENON1T/s2only_data_release}{official response release},
pinned locally in the G2-X provenance card.
\end{thebibliography}
\end{document}
